% ID: SEL-D0031-OB-E01
% ESTADO: cotejada. Segunda pasada visual completa y borrador compilado revisados.
% Fuente: Selectividad/matematicas_aplicadas/titular_junio_6 EXAMEN MAT APLIC C S.pdf
% SHA-256: bebc339cfd8f69c7c4fae160fd165eaef5f6bc93cfd6d15cb9ce289b7029d167
% Página: 2; opción B; ejercicio original 1.

\textbf{(2.5 puntos)} Un taller fabrica y vende dos tipos de alfombras, de seda y de lana. Para la elaboración de una unidad se necesita un trabajo manual de 2 horas para el primer tipo y de 3 horas para el segundo y de un trabajo de máquina de 2 horas para el primer tipo y de 1 hora para el segundo. Por cuestiones laborales y de planificación, se dispone de hasta 600 horas al mes para el trabajo manual y de hasta 480 horas al mes para el destinado a la máquina.

Si el beneficio por unidad para cada tipo de alfombra es de 150 \texteuro{} y 100 \texteuro{}, respectivamente, \textquestiondown{}cuántas alfombras de cada tipo debe elaborar para obtener el máximo beneficio? \textquestiondown{}A cuánto asciende el mismo?
